\documentclass[10pt,a4paper]{amsart}
\usepackage[margin=19mm]{geometry}
\usepackage{amsmath,amssymb,mathtools}
\usepackage[scr=rsfs,cal=euler]{mathalfa}
\usepackage{xfrac}
\usepackage[unicode,hidelinks,bookmarks=false]{hyperref}
\newcommand{\ie}{i.e.\ }
\newcommand{\cf}{cf.\ }
\newcommand{\resp}{respectively\ }
\newcommand{\Subsection}{\subsection}
\providecommand{\nobreakdash}{\nobreak\hbox{-}}
\setlength{\emergencystretch}{2em}
\multlinegap0pt
\usepackage{bm}
\usepackage{bbm}
\usepackage{dsfont}
\usepackage{xspace}
\usepackage{cancel}
\usepackage{mathtools}
\usepackage{amsthm}
\makeatletter
\@ifundefined{theorem}{\newtheorem{theorem}{Theorem}}{}
\makeatother
\title[Some more constructions of $n-$cycle permutation polynomials]{Some more constructions of $n-$cycle permutation polynomials}
\author{Varsha Jarali, Prasanna Poojary,, Vadiraja Bhatta G. R}
\date{Source: arXiv:2506.21936v1 (CC BY 4.0)}
\begin{document}
\maketitle

\noindent\emph{Source and licence.} Some more constructions of $n-$cycle permutation polynomials. Version arXiv:2506.21936v1 (CC BY 4.0), distributed under the Creative Commons Attribution 4.0 International licence, as recorded in the arXiv licence metadata for this exact version and shown on the abstract page https://arxiv.org/abs/2506.21936v1. Full rights notice: licenses/arxiv-2506.21936-CC-BY-4.0.txt.\par\smallskip

\noindent\textit{Editorial scope.} Counted unit: the Theorem whose source label is \texttt{t5} in the article (source0.tex lines 483--513, source label \texttt{t5}), delivered together with the article's own complete original proof of it (source0.tex lines 514--608, one \texttt{proof} environment of 95 lines). No mathematics is written, repaired or re-derived: statement and proof are the article's own text line for line. Required context retained uncounted: none. Every definition, display and label of those lines is retained unchanged. Omitted from this package and not redistributed: 1--482, 609--721. Nothing inside a retained excerpt is omitted or altered: every retained body is delivered exactly as the article's lines are written, so no omission receipt of any kind accompanies this package. Citations: the retained excerpts carry no citation command at all, so no bibliography entry of the article is transcribed and no reference is redistributed. Reference adaptation: none was needed and none was made; every reference inside the delivered unit resolves inside it, so no unresolved reference or citation is produced. Printed slips kept literal: no printed word, symbol, hypothesis or number of the counted statement or of its proof was altered. Layout and class: the article's own class is \texttt{article} with packages including geometry, graphicx, mathtools, amssymb, amsthm, thmtools, multirow, babel; the wrapper is the lane's pinned \texttt{amsart} editorial wrapper at 10pt on A4, which restates the theorem-like environments the retained text needs and adds the article's own macro definitions that the retained text needs (none), with extra wrapper packages (bm, bbm, dsfont, xspace, cancel, mathtools). The delivered numbering is the unit's own. Assets: no retained span contains a figure environment, a graphics-inclusion command, a PDF-page inclusion, a table or a TikZ picture; no image file is copied into this package, the package declares no asset dependency and no image file is redistributed with it. Licence: version arXiv:2506.21936v1 of this article is distributed under the Creative Commons Attribution 4.0 International licence, as recorded in the arXiv licence metadata for this exact version and shown on the abstract page https://arxiv.org/abs/2506.21936v1. A full rights notice is delivered at licenses/arxiv-2506.21936-CC-BY-4.0.txt. Characters: every retained excerpt is pure ASCII, so the delivered engine, XeLaTeX with its default Latin Modern text font, reports no missing character.
\begin{theorem}\label{t5}
	Let $ F(x)=ax^{2^{i}}+bx^{2^{j}} $, $ a\in F_{2^{m}}^{*} $, and $ b\in F_{2^{m}}^{*} $, where $ i<j<m $. Then either of the following will hold;
	\begin{enumerate}
		\item For $ m $ relatively prime with both 2 and 3, $ F(x) $ can never be a triple cycle permutation polynomial.
		\item For $ m=3k $, or $ 2m=3k $, $ F(x) $ is a triple cycle if and only if $ j=i+k $ and
		
		either,
		\begin{equation}
			i=0\hspace{.2cm} and \hspace{.2cm}  a^{2}+b^{2}=1 \hspace{.2cm} and \hspace{.2cm} a^{2}b=0,  
		\end{equation}
		or, 
		\begin{equation}
			3i=2k,\hspace{.2cm}  a^{2^{2k}+1}=b^{2^{2k}+1},  ~a=a^{2^{i}}b^{2^{i}}(a^{2^{i}}+b^{2^{i}}) \hspace{.2cm} and\hspace{.2cm}  a^{2}+ab=1,
		\end{equation}
		
		or, 
		\begin{equation}
			3i=k,\hspace{.2cm}  a^{2^{k}+1}=b^{2^{k}+1}, ~ b=a^{2^{i}}b^{2^{i}}(a^{2^{i}}+b^{2^{i}}) \hspace{.2cm} and\hspace{.2cm}  b^{2}+ab=1. 
		\end{equation}
		\item For $ m=2k $, $ F(x) $ is a triple cycle if and only if $ i=j+k $ and 
		
		either,
		\begin{equation}
			i=0,\hspace{0.1cm}  a\in F_{2^{k}}, ~ b\in F_{2^{k}}, ~a^{2}+ab=1,\hspace{0.1cm} and\hspace{0.1cm}  ab+b^{3}=0,
		\end{equation}
		or,
		\begin{equation}
			j=0,\hspace{.2cm}  a\in F_{2^{k}} , b\in F_{2^{k}} \hspace{.2cm} and\hspace{.2cm}  ab+b^{2}=1.
		\end{equation}
	\end{enumerate}
\end{theorem}

\begin{proof}
	\begin{equation}
		\begin{split}
			F^{2}(x)&=a(ax^{2^{i}}+bx^{2^{j}})^{2^{i}}+b(ax^{2^{i}}+bx^{2^{j}})^{2^{j}}\\
			&=a^{2^{i}+1}x^{2^{2i}}+b^{2^{j}+1}x^{2^{2j}}+x^{2^{(i+j)}}(ab^{2^{i}}+a^{2^{j}}b)
		\end{split}
	\end{equation}
	\begin{equation}
		\begin{split}
			F^{3}(x)&=a^{2^{i}+1}(ax^{2^{i}}+bx^{2^{j}})^{2^{2i}}+b^{2^{j}+1}(ax^{2^{i}}+bx^{2^{j}})^{2^{2j}}+\\
   &(ax^{2^{i}}+bx^{2^{j}})^{2^{(i+j)}}(ab^{2^{i}}+a^{2^{j}}b)\\
			&=a^{2^{3i}+1}x^{2^{3i}}+b^{2^{3j}+1}x^{2^{3j}}+x^{2^{(2i+j)}}(a^{2^{i}+1}b^{2^{2i}}+a^{2^{(i+j)}+1}b^{2^{i}}+a^{2^{(i+2j)}}b)\\
			&+x^{2^{(i+2j)}}(a^{2^{2j}}b^{2^{j}+1}+ab^{2^{(2i+j)}}+a^{2^{j}}b^{2^{(i+j)}+1})	
		\end{split}	
	\end{equation}
	The exponents $ e $ of $ x $  belong to the set $\{3i,3j,2i+j,i+2j\}$. To get $ F^{3}(x)=x $ (noting that $ i\not=j $) it is required to remove three out of the four values from the set, which is explained in the following 6 cases.
	\begin{enumerate}
		\item  $ 3i\equiv3j(\mod m) $ implies $ 3i=3j+m $, which is impossible if $ 3\not|m $. On the other hand if $  m=3k $ then $ i=j+k $.
		\item $ 3i\equiv2i+j(\mod m) $ implies $ 3i=2i+j+m $ and hence $ i=j+m $ which is impossible. 
		\item $ 3i\equiv2j+i(\mod m) $ implies $ 3i=2j+i+m $ and hence $ 2i=2j+m $ which is impossible for $2\not|m$. On the other hand if $ m=2k $ then $ 2i=2j+2k $ implying  $ i=j+k $.
		\item $ 3j\equiv2i+j(\mod m) $ implies $ 3j=2i+j+m $ which is impossible if $2\not|m$. On the other hand if $ m=2k $ then $ 2j=2i+2k $ implying $ j=i+k $.
		\item $ 3j\equiv i+2j(\mod m) $ implies  $ 3j=2j+i+m $ implying $ j=i+m $ which is impossible.
		\item $ 2i+j\equiv i+2j(\mod m) $ implies $ 2i+j=i+2j+m $ and hence $ i=j+m $ which is impossible.
	\end{enumerate}
	When $ i=j+k $ and $ m=3k $ the form of  $ F^{3}(x) $ will be
	\begin{equation}
		\begin{split}
			F^{3}(x)&=(a^{2^{3i}+1}+b^{2^{3j}+1})x^{2^{3i}}+(a^{2^{i}+1}b^{2^{2i}}+a^{2^{i+j}+1}b^{2^{i}}+a^{2^{i+2j}}b)x^{2^{2i+j}}+\\
			&x^{2^{i+2j}}(a^{2^{2j}}b^{2^{j}+1}+ab^{2^{2i+j}}+a^{2^{j}}b^{2^{i+j}+1}).
		\end{split}
	\end{equation}
	The case  $ 3i=m $ possible only when $ i=0 $ and in this case, $ F^{3}(x)=x $ for all $ x $ if and only if $ a^{2}+b^{2^{3k}+1}=a^{2}+b^{2}=1$, $ a^{2}b+a^{2^{k}+1}b+a^{2^{2k}}b=0 $ and $ b^{2^{k}+1}a^{2^{2k}}+a^{2^{k}}b^{2^{k}+1}+ab^{2^{k}}=0 $. This is possible if $ a^{2}+b^{2}=1  $ and $ a^{2}b=0   $. Otherwise, we must have either $ 2i+j=m $ or $ i+2j=m $. If $ 2i+j=m $, we have  $ 3i=2k $ as $ j=k+i $ and in this case, $ F^{3}(x)=x $ for all $ x $ if and only if 
	\begin{equation}
		a^{2^{2k}+1}+b^{2^{2k}+1}=0,\hspace{.1cm} a^{2^{i}+1}b^{2^{2i}}+a^{2^{2i}+1}b^{2^{i}}+ab=1\hspace{.2cm} and 
	\end{equation}
	\begin{equation}
		b^{2^{i}+1}a^{2^{2i}}+a^{2^{i}}b^{2^{2i}+1}+ab=0.
	\end{equation}
	This is possible if $ a^{2^{2m}+1}=b^{2^{2m}+1} $, $ a=a^{2^{i}}b^{2^{i}}(a^{2^{i}}+b^{2^{i}}) $ and $ a^{2}+ab=1 $.
	
	
	Similarly, if $ i+2j=m $, we have  $ 3i=k $ as $ j=i+k $ and in this case, $ F^{3}(x)=x $ for all $ x $ if and only if
	\begin{equation}
		a^{2^{k}+1}+b^{2^{k}+1}=0,\hspace{.1cm} a^{2^{i}+1}b^{2^{2i}}+a^{2^{2i}+1}b^{2^{i}}+ab=0\hspace{.2cm} and 
	\end{equation}
	\begin{equation}
		b^{2^{i}+1}a^{2^{2i}}+a^{2^{i}}b^{2^{2i}+1}+ab^{2^{2m}}=1.
	\end{equation}
	This is possible if $ a^{2^{k}+1}=b^{2^{k}+1} $, $ b=a^{2^{i}}b^{2^{i}}(a^{2^{i}}+b^{2^{i}}) $ and $ b^{2}+ab=1 $.
	
	
	Now it remains to look at the $3^{rd}$ and $4^{th}$ cases. In both these cases  $ F(x) $ is a triple cycle permutation polynomial and when $ i=j+k $ with $ m=2k $, we have 
	\begin{equation}
		\begin{split}
			F^{3}(x)&=(a^{2^{3i}+1}+a^{2^{2j}}b^{2^{j}+1}+ab^{2^{2i+j}}+a^{2^{j}}b^{2^{i+j}+1})x^{2^{3i}}+\\
			&(b^{2^{3i}+1}+a^{2^{i}+1}b^{2^{2i}}+a^{2^{i+j}+1}b^{2^{i}}+a^{2^{i+2j}}b)x^{2^{3j}}.
		\end{split}
	\end{equation}
	The case $ 3i=m $ is possible only when $ i=0 $. 
	
	Consider two terms
	\begin{equation}\label{e1}
		a^{2}+b^{2^{k}+1}a+ab^{2^{k}}+a^{2^{k}}b^{2^{k}+1}\hspace{.2cm}and
	\end{equation}
	\begin{equation}\label{e2}
		b^{2^{k}+2}+a^{2}b+a^{2^{k}+1}b+ab,
	\end{equation}
	and if $ a\in F_{2^{k}} $, and $ b\in F_{2^{k}} $ Equations (\ref{e1}) and (\ref{e2}) reduce to, 
	\begin{equation}
		a(a+b^{2^{k}})+b(ab^{2^{k}}+a^{2^{k}}b^{2^{k}})=a^{2}+ab,
	\end{equation}
	\begin{equation}
		b(b^{2^{k}+1}+a^{2}+a^{2^{k}+1}+a)=ab+b^{3}.
	\end{equation}
	
	In this case, $ F^{3}(x)=x $ for all $ x $ if and only if $ a\in F_{2^{k}} $, $ b\in F_{2^{k}} $, $ a^{2}+ab=1 $, and $ ab+b^{3}=0 $.
	Otherwise, we must have $ 3j=m $ and is only when $ j=0 $.
	
	Consider two terms 
	\begin{equation}\label{e3}
		a^{2^{k}+2}+b^{2}a+ab+ab^{2^{k}+1}\hspace{.2cm} and
	\end{equation}
	\begin{equation}\label{e4}
		b^{2}+a^{2^{k}+1}b+a^{2^{k}+1}b^{2^{k}}+a^{2^{k}}b.
	\end{equation}
	and if $ a\in F_{2^{k}} $, and $ b\in F_{2^{k}} $ Equations (\ref{e3}) and (\ref{e4}) reduce to
	\begin{equation}
		a(a^{2^{k}+1}+b^{2}+b+b^{2^{k}+1})=0,
	\end{equation}
	\begin{equation}
		b^{2}+a^{2}b+a^{2}b+ab=ab+b^{2}.
	\end{equation}
	In this case, $ F^{3}(x)=x $ for all $ x $ if and only if $ a\in F_{2^{k}} $, $ b\in F_{2^{k}} $ and $ ab+b^{2}=1 $.
	Similarly, for  $ 3j\equiv2i+j(\mod m) $, we get same condition for $ F(x) $ to be triple cycle permutation polynomial over $ F_{2^{m}} $.
\end{proof}

\end{document}
